% Partie : sets-functions-relations
% Chapitre : functions
% Section : partial-functions

\documentclass[../../../include/open-logic-section]{subfiles}


\begin{document}

\olfileid[fr]{sfr}{fun}{par}

\olsection{Fonctions partielles}

\begin{explain}
Il est parfois utile d’assouplir la définition d’une fonction en
n’exigeant plus que sa valeur soit définie pour toute entrée possible.
On parle alors de \emph{fonctions partielles}.
\end{explain}

\begin{defn}
Une \emph{fonction partielle} $f\colon A \pto B$ est une correspondance
qui associe à chaque !!{element} de~$A$ au plus un !!{element} de~$B$.
Si $f$ associe un élément de~$B$ à $x \in A$, on dit que $f(x)$ est
\emph{définie} ; sinon, on dit qu’elle est \emph{non définie}.
Si $f(x)$ est définie, on écrit $f(x) \fdefined$ ; sinon,
on écrit $f(x) \fundefined$. Le \emph{domaine} d’une fonction
partielle~$f$ est la partie de~$A$ où elle est définie :
$\dom{f} = \Setabs{x \in A}{f(x) \fdefined}$.
\end{defn}

\begin{ex}
Toute fonction $f\colon A \to B$ est aussi une fonction partielle.
Les fonctions partielles définies partout sur~$A$, c’est-à-dire celles
que nous avons jusqu’ici simplement appelées des fonctions, sont aussi
appelées des fonctions \emph{totales}.
\end{ex}

\begin{ex}
La fonction partielle $f\colon \Real \pto \Real$ donnée par $f(x) = 1/x$
n’est pas définie pour $x = 0$, et elle est définie partout ailleurs.
\end{ex}

\begin{prob}
Étant donnée $f\colon A \pto B$, définissons la fonction partielle
$g\colon B \pto A$ de la manière suivante : pour tout $y \in B$,
s’il existe un unique $x \in A$ tel que $f(x) = y$, on pose $g(y) = x$ ;
sinon, $g(y) \fundefined$. Montrez que, si $f$ est injective, alors
$g(f(x)) = x$ pour tout $x \in \dom{f}$, et $f(g(y)) = y$ pour tout
$y \in \ran{f}$.
\end{prob}

\begin{defn}[Graphe d’une fonction partielle]
Soit $f\colon A \pto B$ une fonction partielle. Le \emph{graphe} de~$f$
est la relation $R_f \subseteq A \times B$ définie par
\[
R_f = \Setabs{\tuple{x,y}}{f(x) = y}.
\]
\end{defn}

\begin{prop}
Supposons que $R \subseteq A \times B$ vérifie la propriété suivante :
chaque fois que $Rxy$ et $Rxy'$, on a $y = y'$.
Alors $R$ est le graphe de la fonction partielle $f\colon A \pto B$
définie ainsi : s’il existe $y$ tel que $Rxy$, on pose $f(x) = y$ ;
sinon, $f(x) \fundefined$. Si $R$ est de plus \emph{sérielle},
c’est-à-dire si, pour chaque $x \in A$, il existe $y \in B$ tel que
$Rxy$, alors $f$ est totale.
\end{prop}

\begin{proof}
Supposons qu’il existe $y$ tel que $Rxy$. S’il existait aussi $y' \neq y$
tel que $Rxy'$, la condition sur $R$ serait violée. Ainsi, un tel $y$,
lorsqu’il existe, est unique ; $f$ est donc bien définie.
On a manifestement $R_f = R$, et $f$ est totale si~$R$ est sérielle.
\end{proof}

\end{document}
